% Generated draft: requires mathematical and source review.
\documentclass[11pt]{article}
\usepackage[margin=25mm]{geometry}
\usepackage{fontspec}
\setmainfont{Latin Modern Roman}
\usepackage{amsmath,amssymb,mathtools,unicode-math}
\setmathfont{Latin Modern Math}
\usepackage{enumitem,needspace,xurl,hyperref}
\setlength{\parindent}{0pt}
\setlength{\parskip}{5pt}
\setlength{\emergencystretch}{3em}
\newcommand{\N}{\mathbb{N}}
\newcommand{\Z}{\mathbb{Z}}
\newcommand{\Q}{\mathbb{Q}}
\newcommand{\R}{\mathbb{R}}
\newcommand{\C}{\mathbb{C}}
\newcommand{\F}{\mathbb{F}}
\newcommand{\A}{\mathbb{A}}
\newcommand{\PP}{\mathbb{P}}
\newcommand{\E}{\mathbb{E}}
\newcommand{\eps}{\varepsilon}
\newcommand{\dd}{\,\mathrm{d}}
\begin{document}
\section*{UnsolvedMath ID: 30003298}
\textbf{Source problem code:} catalogue entry
\subsection{Definitions and mathematical statement}
The supplied corrected formulation concerns a torsion-free subgroup \(\Gamma\leq\operatorname{Mod}_g\) of finite index and the curve complex \(\mathcal C_g\). The record does not specify the range of the genus \(g\), the surface conventions, or the coefficient ring.

Under the usual closed-surface interpretation, let \(S_g\) be a closed, connected, oriented surface of genus \(g\), and let \(\operatorname{Mod}_g\) be its group of orientation-preserving homeomorphisms modulo isotopy. For \(g\geq 2\), the curve complex \(\mathcal C_g\) has vertices given by isotopy classes of essential simple closed curves on \(S_g\); a finite set of distinct vertices spans a simplex when its classes admit pairwise disjoint representatives. The mapping class group acts on this complex and hence on its homology. These are conventional interpretations, not additional hypotheses verified from the supplied report.

With integral coefficients as an explicit provisional convention, put
\[
M_g=H_{2g-2}(\mathcal C_g;\Z),
\]
viewed as a \(\Z\Gamma\)-module by restricting the induced mapping class group action. The corrected question, under these conventions, is:
\[
\text{For every torsion-free finite-index subgroup }\Gamma\leq\operatorname{Mod}_g,
\quad
H^{2g-1}(\Gamma;M_g)
\text{ is not finitely generated as an abelian group?}
\]
Here \(H^k(\Gamma;M)\) denotes group cohomology with coefficients in the \(\Gamma\)-module \(M\). The intended genus range must be supplied before this can be read as a fully quantified conjecture; the coefficient and finite-generation conventions also require confirmation.

There is a material conflict within the record. Its original formulation calls \(\mathcal C_g\) the universal curve and asks about
\[
H^{2g-1}\bigl(\Gamma;H^{2g-2}(\mathcal C_g)\bigr),
\]
whereas its statement and clean statement call \(\mathcal C_g\) the curve complex and use homology \(H_{2g-2}\). These formulations must not be identified. The draft above follows the supplied corrected formulation provisionally; the original universal curve, its base, and its coefficient action are not specified sufficiently to formalize that alternative.
\subsection{Short English statement}
For a torsion-free finite-index subgroup of the genus-\(g\) mapping class group, the supplied corrected question asks whether its degree-\(2g-1\) group cohomology, with coefficients in the degree-\(2g-2\) homology of the curve complex, is infinitely generated. The genus range and coefficient conventions are missing, and the original formulation uses a different space and cohomology instead of homology. The supplied open status requires independent review.
\subsection{Sources}
\begin{itemize}
\item Supplied local UnsolvedMath registry record, permanent ID \(30003298\), title: Infinite Generation in Mapping-Class-Group Cohomology; exact source problem code: \texttt{catalogue entry}. The record belongs to the Oberwolfach Reports Open Problems collection and supplies proposed year 2016; no individual proposer is named.
\item Supplied report citation: \emph{Surface Bundles} (2017), Oberwolfach Reports, pp.\ 3149--3195, DOI \texttt{10.4171/owr/2016/56}. Supplied source URL: \url{https://doi.org/10.4171/owr/2016/56}.
\item Supplied literature source, labeled OWR 56 (2017): \url{https://doi.org/10.4171/owr/2016/56}. The literature assessment dated \texttt{2026-08-21} labels the problem open and reports that no checked primary source proves or disproves infinite generation of the stated mapping-class-group cohomology group. This is supplied dated curation derived from the OWR record, not an independently verified literature assessment; time-sensitive status and cited claims require independent verification.
\item The registry marks the statement \texttt{corrected\_verified} and reports that it was checked against the official Oberwolfach report. That supplied verification claim has not been independently checked here, and the conflict with the original formulation remains recorded.
\end{itemize}
\end{document}
